\begin{pblproject}
\label{sec:05-rigor-check:06-checkpoint-project:1}
\textbf{\large Checkpoint project: a Monoid from scratch}\par\smallskip


Part I is now complete: terms and types, structures, propositions and proofs, tactics, and the rigor questions just answered. Before Chapter 6 builds \texttt{Group}, this project asks for a smaller, self-contained version of exactly that construction --- a \textbf{monoid}: a set with an associative operation and a two-sided identity, but \emph{no} inverses. Everything needed is already in hand from Chapters 1--5 alone; nothing here depends on \texttt{Group}, so it previews the next chapter's construction rather than spoiling it.

\textbf{Learning objectives.} Practice bundling data with proof obligations in a \texttt{structure} (Chapter 2), writing tactic-mode proofs for a genuine algebraic axiom set (Chapter 4), and building a small ``prove it once, generically'' theorem in the style Chapter 7 will use for \texttt{Group} --- here done for the weaker \texttt{Monoid} first.

\textbf{Prerequisites.} Chapters 1--5 only.

\textbf{Milestones.}

\begin{enumerate}
\def\labelenumi{\arabic{enumi}.}
\tightlist
\item
  Define \texttt{structure Monoid (M : Type) where} with fields \texttt{op}, \texttt{id}, and proof obligations \texttt{assoc}, \texttt{id\_\allowbreak{}left}, \texttt{id\_\allowbreak{}right} --- exactly \texttt{Group}'s fields (Chapter 6 §2), minus \texttt{inv}/\texttt{inv\_\allowbreak{}left}/\texttt{inv\_\allowbreak{}right}.
\item
  Build at least one concrete instance. Two natural choices, both reusing facts already available from core Lean: \texttt{List α} under \texttt{++}/\texttt{[]}, or \texttt{Nat} under \texttt{*}/\texttt{1}.
\item
  Prove a generic theorem for an arbitrary \texttt{Mn : Monoid M}: \emph{the identity is unique} --- if \texttt{e' : M} satisfies \texttt{∀ a, Mn.op e' a = a}, then \texttt{e' = Mn.id}. This is exactly Chapter 7's Theorem 1 (\texttt{id\_\allowbreak{}unique}) carried out one chapter early, and --- worth checking directly --- its proof never needs an inverse, so nothing about lacking one gets in the way.
\end{enumerate}

\textbf{Deliverable.} \texttt{structure Monoid}, at least one concrete instance (with all three proof obligations discharged), and the generic uniqueness theorem, applied to that instance.

\textbf{Self-verification.}

\begin{lstlisting}
structure Monoid (M : Type) where
  op : M (*$\rightarrow$*) M (*$\rightarrow$*) M
  id : M
  assoc : (*$\forall$*) a b c : M, op (op a b) c = op a (op b c)
  id_left : (*$\forall$*) a : M, op id a = a
  id_right : (*$\forall$*) a : M, op a id = a

def listMonoid ((*$\alpha$*) : Type) : Monoid (List (*$\alpha$*)) where
  op := List.append
  id := []
  assoc := by intro a b c; exact List.append_assoc a b c
  id_left := by intro a; exact List.nil_append a
  id_right := by intro a; exact List.append_nil a

theorem monoid_id_unique {M : Type} (Mn : Monoid M) (e' : M)
    (h : (*$\forall$*) a : M, Mn.op e' a = a) : e' = Mn.id := by
  have step1 : Mn.op e' Mn.id = Mn.id := h Mn.id
  have step2 : Mn.op e' Mn.id = e' := Mn.id_right e'
  rw [(*$\leftarrow$*) step2]
  exact step1

-- Applying the generic theorem to the concrete instance costs nothing
-- beyond naming it (*\textemdash{}*) the same "prove once, use everywhere" payoff
-- Chapter 6 (*\S*)6 promises for Group, delivered here one chapter early.
#check monoid_id_unique (listMonoid Nat) [] (fun a => List.nil_append a)
\end{lstlisting}

If this compiles (\texttt{lake env lean} on a file containing it, or pasted into \texttt{lean\_\allowbreak{}project}), the project is done. A full worked solution, including a second instance, is in \hyperref[sec:14-appendix-solutions:04-chapter-5:1]{Appendix, Chapter 5}.

\end{pblproject}
